\chapter{Obiectivele proiectului}\label{ch:obiective}
\pagestyle{fancy}

Proiectul de față își propune realizarea unui sistem de control pe FPGA pentru o
mână robotică acționată de 8 servomotoare. Sistemul primește, printr-o interfață
serială, unghiurile articulațiilor unei mâini umane urmărite de un sistem extern
de viziune și le transformă în comenzi de poziție pentru servomotoare. Pentru ca
mișcarea reprodusă să fie naturală și lipsită de tremur, datele primite sunt
prelucrate de un filtru Kalman înainte de a fi trimise spre actuatoare. Pe lângă
modul de operare comandat extern, sistemul poate funcționa și autonom, parcurgând
secvențe de mișcare predefinite, iar separarea dintre logica de control și partea
de putere este asigurată prin izolare galvanică.

\section{Preluarea și filtrarea datelor pentru mimarea fidelă a mișcării}
\label{sec:obj-principal}

Obiectivul principal al proiectului este preluarea unghiurilor articulațiilor
printr-o interfață serială și prelucrarea acestora astfel încât mâna robotică să
reproducă cât mai fidel mișcarea mâinii umane. Datele de intrare descriu, pentru
fiecare articulație, unghiul dorit, iar sistemul trebuie să convertească aceste
valori în comenzi de poziție stabile, transmise în mod continuu către
servomotoare. Calitatea acestei reproduceri reprezintă criteriul central după care
este evaluat întregul sistem.

Datele primite de la sistemul de viziune nu sunt însă perfecte, deoarece conțin
zgomot și variații de la un cadru la altul, cauzate de erorile inerente ale
procesului de detecție. Dacă aceste valori ar fi trimise direct la servomotoare, mișcarea
rezultată ar fi sacadată, cu tremurături vizibile și nepotrivite cu gestul real.
Pentru a corecta acest neajuns, sistemul aplică un filtru Kalman, care estimează
poziția reală a fiecărei articulații pe baza valorilor zgomotoase primite și
produce o traiectorie netedă și continuă. Filtrarea Kalman constituie astfel
elementul central al lucrării, fiind componenta care transformă un flux brut de
date într-o mișcare fluidă, apropiată de cea naturală.


\section{Acționarea servomotoarelor pentru manipularea obiectelor}
\label{sec:obj-pwm}

Un obiectiv secundar al proiectului este acționarea precisă a celor 8
servomotoare, astfel încât mâna robotică să poată executa gesturi utile, inclusiv
apucarea și manipularea unor obiecte. Fiecare servomotor este controlat printr-un
semnal de tip PWM, a cărui durată (\textit{pulse width}) determină poziția
unghiulară a articulației corespunzătoare. Pentru a obține o mișcare coordonată a degetelor,
toate cele 8 semnale trebuie generate simultan și actualizate sincron, fără
întârzieri între canale. Tocmai aici intervine avantajul platformei FPGA. Spre
deosebire de un microcontroler, care ar genera semnalele secvențial, FPGA-ul
permite implementarea a 8 module PWM independente, ce funcționează în paralel și
produc cele 8 semnale în mod perfect simultan.

Acuratețea și stabilitatea acestor semnale sunt esențiale deoarece o variație nedorită a
duratei pulsului se traduce printr-o tremurătură a articulației, iar lipsa
sincronizării între canale ar face ca degetele să se miște la momente ușor
diferite, compromițând prinderea fermă a unui obiect. Întrucât în FPGA fiecare
semnal este generat de un circuit dedicat, sincronizat cu ceasul de sistem,
acestea sunt precise și independente, îndeplinind cerințele acestui obiectiv.

Pentru manipularea efectivă a obiectelor, nu este suficient ca fiecare deget să
ajungă într-o poziție dorită, ci este nevoie ca toate degetele să își coordoneze
mișcarea într-un gest unitar. Apucarea unui obiect presupune ca degetele să se
închidă progresiv și împreună, până când obiectul este prins ferm. Dacă un singur
canal ar fi întârziat sau instabil, prinderea ar deveni nesigură, iar obiectul ar
putea aluneca. Din acest motiv, generarea simultană și sincronă a tuturor celor 8
semnale nu este doar o cerință tehnică, ci condiția fundamentală pentru ca mâna
robotică să poată îndeplini sarcini reale de manipulare.

\section{Funcționarea autonomă}
\label{sec:obj-autonom}

Un alt obiectiv secundar este capacitatea sistemului de a funcționa autonom, fără
a depinde de conexiunea cu un calculator extern. În acest mod, mâna robotică
parcurge o secvență de mișcări predefinite, stocate intern, permițând
demonstrarea funcționalității sistemului în lipsa sursei de date externe.

Acest mod de operare este util atât pentru testarea independentă a părții de
control și actuare, cât și pentru prezentarea sistemului în condiții în care
echipamentul de captură nu este disponibil. Comutarea între modul comandat extern
și cel autonom trebuie să fie simplă și să nu necesite reconfigurarea sistemului.

Secvențele predefinite urmăresc demonstrarea principalelor capabilități ale mâinii
robotice, printre care deschiderea și închiderea completă a degetelor, mișcările
individuale ale fiecărui deget, precum și gesturi uzuale. În acest fel, modul autonom oferă o
imagine clară asupra întregului interval de mișcare al sistemului și asupra
funcționării corecte a fiecărui servomotor în parte, fără a depinde de o sursă
externă de date.

Din punct de vedere practic, acest obiectiv prezintă și un avantaj important pe
parcursul dezvoltării. Întrucât partea de viziune și partea de control sunt
realizate separat, modul autonom permite verificarea completă a sistemului de
control înainte de integrarea cu sistemul de captură, eliminând dependența de
disponibilitatea acestuia din urmă. Astfel, eventualele probleme de actuare pot fi
identificate și corectate independent, ceea ce simplifică procesul de testare și
validare.

\section{Izolarea galvanică și siguranța}
\label{sec:obj-izolare}

Deoarece sistemul cuprinde două domenii electrice distincte, partea de logică de
control la tensiune redusă și partea de putere care alimentează servomotoarele,
un obiectiv secundar important este asigurarea unei separări galvanice între
acestea. Această separare protejează componenta digitală de eventualele
perturbații, vârfuri de curent sau zgomot generate de servomotoare în timpul
funcționării.

Necesitatea acestei separări devine evidentă dacă se ține cont de comportamentul
electric al servomotoarelor. La pornire și în timpul mișcării, acestea pot
genera vârfuri de curent și perturbații pe linia de alimentare, fenomene care, în
absența unei separări, s-ar putea propaga către placa de control și ar putea
afecta funcționarea logicii digitale sau chiar deteriora componentele acesteia.

Izolarea galvanică asigură atât siguranța în funcționare, cât și integritatea
semnalelor transmise între cele două domenii. În acest scop, semnalele de comandă
către servomotoare și cele de comunicație sunt transferate prin componente
dedicate de izolare, care permit trecerea informației fără a exista o legătură
electrică directă între cele două părți. Îndeplinirea acestui obiectiv contribuie
la fiabilitatea generală a sistemului și la protejarea plăcii de control împotriva
defectelor care ar putea proveni din partea de putere.

\section{Arhitectură modulară și încadrarea în resursele disponibile}
\label{sec:obj-modular}

Un ultim obiectiv secundar vizează organizarea sistemului într-o arhitectură
modulară, în care fiecare funcție (comunicație, filtrare, generare a semnalelor,
control) este implementată ca un bloc independent. O astfel de structură ușurează
dezvoltarea, testarea și eventuala extindere ulterioară a sistemului, fiecare
modul putând fi verificat separat.

În plus, întreaga implementare trebuie să se încadreze în resursele limitate ale
plăcii FPGA utilizate și să respecte constrângerile temporale impuse de
funcționarea în timp real. Respectarea acestor limite reprezintă o cerință de
proiectare permanentă, care influențează deciziile luate pe parcursul realizării
fiecărui modul.
